% ------------------------------------------------------------
% A7.1 Collapse-Surface Identification & Classification (v0.1)
% Author: Borealis S. Hedling
% Altitude: A7.1 (Non-Dual Diagnostic Layer)
% Purpose: Define the identification and classification operators
%          required before A7.2 repair can be safely applied.
% ------------------------------------------------------------

\documentclass[11pt]{article}
\usepackage{amsmath, amssymb, amsthm}

\title{A7.1 Collapse-Surface Identification and Classification}
\author{Borealis S. Hedling}
\date{October 2026}

\begin{document}
\maketitle

\section*{Overview}

This document defines the A7.1 diagnostic layer responsible for identifying
and classifying collapse surfaces produced by A7.0 non-dual projection
operators. A7.1 operators do not repair collapse surfaces; they provide the
structural and spectral information required for A7.2 repair.

% ------------------------------------------------------------
\section{Collapse-Surface Identification}

Define the collapse-surface indicator:
\[
\Xi(p) =
\begin{cases}
1 & p \in \Sigma,\\
0 & p \notin \Sigma.
\end{cases}
\]

A7.1 must identify:
\begin{itemize}
  \item curvature singularities,
  \item spectral tears,
  \item modal discontinuities,
  \item holonomy fractures.
\end{itemize}

Identification is non-dual:
\[
\Xi : M \to \{0,1\},
\]
and may be discontinuous.

% ------------------------------------------------------------
\section{Classification Operators}

Define the classification operator:
\[
\mathcal{C}_{\mathrm{class}} : \Sigma \to \mathcal{T},
\]
where $\mathcal{T}$ is the set of collapse types.

\subsection{Curvature Singularity Classification}
\[
\mathcal{C}_{\mathrm{curv}}(p) =
\begin{cases}
\text{Type I} & |K(p)| \to \infty,\\
\text{Type II} & \nabla K(p) \to \infty.
\end{cases}
\]

\subsection{Spectral Tear Classification}
\[
\mathcal{C}_{\mathrm{spec}}(p) =
\delta_{\mathrm{spec}}(p) > \epsilon_{\mathrm{tear}}.
\]

\subsection{Modal Discontinuity Classification}
\[
\mathcal{C}_{\mathrm{modal}}(p) :
\mu(p^-) \not\leadsto \mu(p^+).
\]

\subsection{Holonomy Fracture Classification}
\[
\mathcal{C}_{\mathrm{hol}}(p) :
H(\gamma) > \Theta_{\max}.
\]

% ------------------------------------------------------------
\section{Collapse-Surface Typology}

Define the collapse typology:
\[
\mathcal{T} =
\{\text{Curvature}, \text{Spectral}, \text{Modal}, \text{Holonomy}\}.
\]

Each collapse surface point is assigned a type:
\[
\mathcal{T}(p) =
\mathcal{C}_{\mathrm{class}}(p).
\]

Mixed-type collapse surfaces satisfy:
\[
\mathcal{T}(p) = \{\tau_1, \tau_2, \dots\}.
\]

% ------------------------------------------------------------
\section{Collapse-Surface Geometry}

Define the local collapse geometry:
\[
\Gamma(p) = (K(p), \mathrm{Spec}(p), \mu(p), H(p)).
\]

A7.1 records geometry but does not modify it:
\[
\Gamma_{\mathrm{A7.1}}(p) = \Gamma_{\mathrm{A7.0}}(p).
\]

% ------------------------------------------------------------
\section{Repair Readiness Conditions}

A collapse surface is ready for A7.2 repair if:
\[
\Xi(p) = 1,
\]
\[
\mathcal{T}(p) \text{ is defined},
\]
\[
\Gamma(p) \text{ is finite}.
\]

These conditions ensure A7.2 receives complete diagnostic data.

% ------------------------------------------------------------
\section{Conclusion}

A7.1 defines the identification and classification operators required for
collapse-surface repair. It provides the diagnostic structure needed by
A7.2 and must be executed after A7.0 projection but before any repair
operators are applied.

\end{document}
